% ===========================================================================
% 3. Testbed Design
% ===========================================================================
\section{Testbed Design}
\label{sec:testbed}

\subsection{Architecture Overview}
\label{sec:testbed-arch}

\textsc{agent-traps-lab} is a modular TypeScript framework built on
the Bun runtime~\cite{bun2024}. \Cref{fig:architecture} shows the
high-level architecture. The testbed comprises five principal
components:

\begin{enumerate}[nosep,leftmargin=*]
  \item \textbf{Trap Registry} (\texttt{src/traps/registry.ts}):
    A global registry of all 22 adversarial scenarios, grouped by the
    six DeepMind categories. Each scenario implements the
    \texttt{TrapScenario} interface, which defines four lifecycle
    methods: \texttt{setup()}, \texttt{execute()}, \texttt{evaluate()},
    and \texttt{teardown()}.

  \item \textbf{Experiment Harness} (\texttt{src/harness/}):
    The matrix generator (\texttt{matrix.ts}) produces the full
    factorial experiment design. The runner (\texttt{runner.ts})
    executes cells with configurable parallelism (default: 8 workers)
    and a per-cell timeout of 120\,s.

  \item \textbf{Agent Handle} (\texttt{src/agents/agent-handle.ts}):
    An abstraction over four LLM providers (OpenAI, Anthropic, Google)
    that exposes a uniform \texttt{sendTask()} / \texttt{getDecision()}
    API. Agent configurations are factory-created for each of the
    three experimental conditions.

  \item \textbf{Mitigation Suite} (\texttt{src/mitigations/}):
    Seven defense modules, each implementing the \texttt{Mitigation}
    interface with \texttt{preProcess()} and \texttt{postProcess()}
    hooks. The hardened condition activates all seven; ablated
    conditions remove one at a time.

  \item \textbf{Reporter} (\texttt{src/harness/reporter.ts}):
    Transforms raw experiment results into paper-ready \LaTeX\ assets
    including \texttt{booktabs} tables and \texttt{pgfplots} figures.
\end{enumerate}

\begin{figure}[t]
  \centering
  % ===========================================================================
% Architecture Diagram — Testbed Pipeline
% ===========================================================================
% Included by: % ===========================================================================
% Architecture Diagram — Testbed Pipeline
% ===========================================================================
% Included by: % ===========================================================================
% Architecture Diagram — Testbed Pipeline
% ===========================================================================
% Included by: \input{figures/architecture}
% Shows: datasets → trap scenarios → harness runner → agent handle → metrics → reporter
%
\begin{tikzpicture}[
  % Node styles
  component/.style={
    rectangle,
    draw=black!70,
    fill=blue!8,
    rounded corners=3pt,
    minimum height=1.0cm,
    minimum width=2.2cm,
    font=\small\sffamily,
    align=center,
    thick,
  },
  agent/.style={
    rectangle,
    draw=black!70,
    fill=green!10,
    rounded corners=3pt,
    minimum height=0.75cm,
    minimum width=1.6cm,
    font=\scriptsize\sffamily,
    align=center,
    thick,
  },
  mitigation/.style={
    rectangle,
    draw=black!70,
    fill=orange!10,
    rounded corners=3pt,
    minimum height=0.75cm,
    minimum width=1.6cm,
    font=\scriptsize\sffamily,
    align=center,
    thick,
  },
  data/.style={
    cylinder,
    draw=black!70,
    fill=yellow!10,
    shape border rotate=90,
    aspect=0.25,
    minimum height=0.9cm,
    minimum width=1.4cm,
    font=\small\sffamily,
    align=center,
    thick,
  },
  output/.style={
    rectangle,
    draw=black!70,
    fill=red!8,
    rounded corners=3pt,
    minimum height=0.9cm,
    minimum width=1.8cm,
    font=\small\sffamily,
    align=center,
    thick,
  },
  arrow/.style={
    -Stealth,
    thick,
    black!70,
  },
  darrow/.style={
    Stealth-Stealth,
    thick,
    black!50,
    dashed,
  },
  grouplabel/.style={
    font=\scriptsize\sffamily\itshape,
    text=black!60,
  },
]

% ── Row 1: Data sources ──
\node[data] (datasets) at (0, 0) {Datasets};

% ── Row 2: Trap Registry ──
\node[component] (registry) at (0, -1.8) {Trap Registry\\{\scriptsize 22 scenarios}};

% Category labels under registry
\node[grouplabel, below=0.15cm of registry] (cats) {
  CI(4)~~SM(4)~~CS(4)~~BC(4)~~SY(3)~~HL(3)
};

% ── Row 3: Harness Runner ──
\node[component] (harness) at (0, -4.0) {Harness\\Runner};
\node[component, left=1.0cm of harness] (matrix) {Matrix\\Generator};

% ── Row 4: Agent Handles (4 models) ──
\node[agent] (gpt4o) at (-3.2, -6.0) {GPT-4o};
\node[agent] (claude) at (-1.1, -6.0) {Claude\\Sonnet~4};
\node[agent] (gemini) at (1.0, -6.0) {Gemini\\2.5 Pro};
\node[agent] (mini) at (3.2, -6.0) {GPT-4o\\mini};

% ── Mitigations (floating right of agents) ──
\node[mitigation] (mits) at (5.2, -6.0) {Mitigation\\Suite (7)};

% ── Row 5: Metrics ──
\node[component] (metrics) at (0, -8.0) {Metrics\\Engine};

% ── Row 6: Reporter ──
\node[output] (reporter) at (-2.0, -9.8) {Reporter\\{\scriptsize \LaTeX\ assets}};
\node[output] (results) at (2.0, -9.8) {Results\\{\scriptsize JSON}};

% ── Arrows ──
% Data → Registry
\draw[arrow] (datasets) -- (registry);

% Registry → Harness
\draw[arrow] (registry) -- (harness)
  node[midway, right, font=\scriptsize] {scenarios};

% Matrix → Harness
\draw[arrow] (matrix) -- (harness)
  node[midway, above, font=\scriptsize] {cells};

% Harness → Agent handles
\draw[arrow] (harness) -- (gpt4o)
  node[midway, left, font=\scriptsize, xshift=-2pt] {};
\draw[arrow] (harness) -- (claude);
\draw[arrow] (harness) -- (gemini);
\draw[arrow] (harness) -- (mini);

% Mitigations ↔ Agents
\draw[darrow] (mits) -- (mini)
  node[midway, above, font=\scriptsize] {};
\draw[darrow] (mits.south west) -- (gemini.east);

% Agent handles → Metrics
\draw[arrow] (gpt4o) -- (metrics);
\draw[arrow] (claude) -- (metrics);
\draw[arrow] (gemini) -- (metrics);
\draw[arrow] (mini) -- (metrics);

% Metrics → Outputs
\draw[arrow] (metrics) -- (reporter)
  node[midway, left, font=\scriptsize] {aggregates};
\draw[arrow] (metrics) -- (results)
  node[midway, right, font=\scriptsize] {raw data};

% ── Background boxes ──
\begin{scope}[on background layer]
  % Agent group box
  \node[
    draw=green!40,
    fill=green!3,
    rounded corners=5pt,
    fit=(gpt4o)(claude)(gemini)(mini),
    inner sep=6pt,
    label={[grouplabel]above:{Agent Handles (3 conditions each)}},
  ] {};
\end{scope}

\end{tikzpicture}

% Shows: datasets → trap scenarios → harness runner → agent handle → metrics → reporter
%
\begin{tikzpicture}[
  % Node styles
  component/.style={
    rectangle,
    draw=black!70,
    fill=blue!8,
    rounded corners=3pt,
    minimum height=1.0cm,
    minimum width=2.2cm,
    font=\small\sffamily,
    align=center,
    thick,
  },
  agent/.style={
    rectangle,
    draw=black!70,
    fill=green!10,
    rounded corners=3pt,
    minimum height=0.75cm,
    minimum width=1.6cm,
    font=\scriptsize\sffamily,
    align=center,
    thick,
  },
  mitigation/.style={
    rectangle,
    draw=black!70,
    fill=orange!10,
    rounded corners=3pt,
    minimum height=0.75cm,
    minimum width=1.6cm,
    font=\scriptsize\sffamily,
    align=center,
    thick,
  },
  data/.style={
    cylinder,
    draw=black!70,
    fill=yellow!10,
    shape border rotate=90,
    aspect=0.25,
    minimum height=0.9cm,
    minimum width=1.4cm,
    font=\small\sffamily,
    align=center,
    thick,
  },
  output/.style={
    rectangle,
    draw=black!70,
    fill=red!8,
    rounded corners=3pt,
    minimum height=0.9cm,
    minimum width=1.8cm,
    font=\small\sffamily,
    align=center,
    thick,
  },
  arrow/.style={
    -Stealth,
    thick,
    black!70,
  },
  darrow/.style={
    Stealth-Stealth,
    thick,
    black!50,
    dashed,
  },
  grouplabel/.style={
    font=\scriptsize\sffamily\itshape,
    text=black!60,
  },
]

% ── Row 1: Data sources ──
\node[data] (datasets) at (0, 0) {Datasets};

% ── Row 2: Trap Registry ──
\node[component] (registry) at (0, -1.8) {Trap Registry\\{\scriptsize 22 scenarios}};

% Category labels under registry
\node[grouplabel, below=0.15cm of registry] (cats) {
  CI(4)~~SM(4)~~CS(4)~~BC(4)~~SY(3)~~HL(3)
};

% ── Row 3: Harness Runner ──
\node[component] (harness) at (0, -4.0) {Harness\\Runner};
\node[component, left=1.0cm of harness] (matrix) {Matrix\\Generator};

% ── Row 4: Agent Handles (4 models) ──
\node[agent] (gpt4o) at (-3.2, -6.0) {GPT-4o};
\node[agent] (claude) at (-1.1, -6.0) {Claude\\Sonnet~4};
\node[agent] (gemini) at (1.0, -6.0) {Gemini\\2.5 Pro};
\node[agent] (mini) at (3.2, -6.0) {GPT-4o\\mini};

% ── Mitigations (floating right of agents) ──
\node[mitigation] (mits) at (5.2, -6.0) {Mitigation\\Suite (7)};

% ── Row 5: Metrics ──
\node[component] (metrics) at (0, -8.0) {Metrics\\Engine};

% ── Row 6: Reporter ──
\node[output] (reporter) at (-2.0, -9.8) {Reporter\\{\scriptsize \LaTeX\ assets}};
\node[output] (results) at (2.0, -9.8) {Results\\{\scriptsize JSON}};

% ── Arrows ──
% Data → Registry
\draw[arrow] (datasets) -- (registry);

% Registry → Harness
\draw[arrow] (registry) -- (harness)
  node[midway, right, font=\scriptsize] {scenarios};

% Matrix → Harness
\draw[arrow] (matrix) -- (harness)
  node[midway, above, font=\scriptsize] {cells};

% Harness → Agent handles
\draw[arrow] (harness) -- (gpt4o)
  node[midway, left, font=\scriptsize, xshift=-2pt] {};
\draw[arrow] (harness) -- (claude);
\draw[arrow] (harness) -- (gemini);
\draw[arrow] (harness) -- (mini);

% Mitigations ↔ Agents
\draw[darrow] (mits) -- (mini)
  node[midway, above, font=\scriptsize] {};
\draw[darrow] (mits.south west) -- (gemini.east);

% Agent handles → Metrics
\draw[arrow] (gpt4o) -- (metrics);
\draw[arrow] (claude) -- (metrics);
\draw[arrow] (gemini) -- (metrics);
\draw[arrow] (mini) -- (metrics);

% Metrics → Outputs
\draw[arrow] (metrics) -- (reporter)
  node[midway, left, font=\scriptsize] {aggregates};
\draw[arrow] (metrics) -- (results)
  node[midway, right, font=\scriptsize] {raw data};

% ── Background boxes ──
\begin{scope}[on background layer]
  % Agent group box
  \node[
    draw=green!40,
    fill=green!3,
    rounded corners=5pt,
    fit=(gpt4o)(claude)(gemini)(mini),
    inner sep=6pt,
    label={[grouplabel]above:{Agent Handles (3 conditions each)}},
  ] {};
\end{scope}

\end{tikzpicture}

% Shows: datasets → trap scenarios → harness runner → agent handle → metrics → reporter
%
\begin{tikzpicture}[
  % Node styles
  component/.style={
    rectangle,
    draw=black!70,
    fill=blue!8,
    rounded corners=3pt,
    minimum height=1.0cm,
    minimum width=2.2cm,
    font=\small\sffamily,
    align=center,
    thick,
  },
  agent/.style={
    rectangle,
    draw=black!70,
    fill=green!10,
    rounded corners=3pt,
    minimum height=0.75cm,
    minimum width=1.6cm,
    font=\scriptsize\sffamily,
    align=center,
    thick,
  },
  mitigation/.style={
    rectangle,
    draw=black!70,
    fill=orange!10,
    rounded corners=3pt,
    minimum height=0.75cm,
    minimum width=1.6cm,
    font=\scriptsize\sffamily,
    align=center,
    thick,
  },
  data/.style={
    cylinder,
    draw=black!70,
    fill=yellow!10,
    shape border rotate=90,
    aspect=0.25,
    minimum height=0.9cm,
    minimum width=1.4cm,
    font=\small\sffamily,
    align=center,
    thick,
  },
  output/.style={
    rectangle,
    draw=black!70,
    fill=red!8,
    rounded corners=3pt,
    minimum height=0.9cm,
    minimum width=1.8cm,
    font=\small\sffamily,
    align=center,
    thick,
  },
  arrow/.style={
    -Stealth,
    thick,
    black!70,
  },
  darrow/.style={
    Stealth-Stealth,
    thick,
    black!50,
    dashed,
  },
  grouplabel/.style={
    font=\scriptsize\sffamily\itshape,
    text=black!60,
  },
]

% ── Row 1: Data sources ──
\node[data] (datasets) at (0, 0) {Datasets};

% ── Row 2: Trap Registry ──
\node[component] (registry) at (0, -1.8) {Trap Registry\\{\scriptsize 22 scenarios}};

% Category labels under registry
\node[grouplabel, below=0.15cm of registry] (cats) {
  CI(4)~~SM(4)~~CS(4)~~BC(4)~~SY(3)~~HL(3)
};

% ── Row 3: Harness Runner ──
\node[component] (harness) at (0, -4.0) {Harness\\Runner};
\node[component, left=1.0cm of harness] (matrix) {Matrix\\Generator};

% ── Row 4: Agent Handles (4 models) ──
\node[agent] (gpt4o) at (-3.2, -6.0) {GPT-4o};
\node[agent] (claude) at (-1.1, -6.0) {Claude\\Sonnet~4};
\node[agent] (gemini) at (1.0, -6.0) {Gemini\\2.5 Pro};
\node[agent] (mini) at (3.2, -6.0) {GPT-4o\\mini};

% ── Mitigations (floating right of agents) ──
\node[mitigation] (mits) at (5.2, -6.0) {Mitigation\\Suite (7)};

% ── Row 5: Metrics ──
\node[component] (metrics) at (0, -8.0) {Metrics\\Engine};

% ── Row 6: Reporter ──
\node[output] (reporter) at (-2.0, -9.8) {Reporter\\{\scriptsize \LaTeX\ assets}};
\node[output] (results) at (2.0, -9.8) {Results\\{\scriptsize JSON}};

% ── Arrows ──
% Data → Registry
\draw[arrow] (datasets) -- (registry);

% Registry → Harness
\draw[arrow] (registry) -- (harness)
  node[midway, right, font=\scriptsize] {scenarios};

% Matrix → Harness
\draw[arrow] (matrix) -- (harness)
  node[midway, above, font=\scriptsize] {cells};

% Harness → Agent handles
\draw[arrow] (harness) -- (gpt4o)
  node[midway, left, font=\scriptsize, xshift=-2pt] {};
\draw[arrow] (harness) -- (claude);
\draw[arrow] (harness) -- (gemini);
\draw[arrow] (harness) -- (mini);

% Mitigations ↔ Agents
\draw[darrow] (mits) -- (mini)
  node[midway, above, font=\scriptsize] {};
\draw[darrow] (mits.south west) -- (gemini.east);

% Agent handles → Metrics
\draw[arrow] (gpt4o) -- (metrics);
\draw[arrow] (claude) -- (metrics);
\draw[arrow] (gemini) -- (metrics);
\draw[arrow] (mini) -- (metrics);

% Metrics → Outputs
\draw[arrow] (metrics) -- (reporter)
  node[midway, left, font=\scriptsize] {aggregates};
\draw[arrow] (metrics) -- (results)
  node[midway, right, font=\scriptsize] {raw data};

% ── Background boxes ──
\begin{scope}[on background layer]
  % Agent group box
  \node[
    draw=green!40,
    fill=green!3,
    rounded corners=5pt,
    fit=(gpt4o)(claude)(gemini)(mini),
    inner sep=6pt,
    label={[grouplabel]above:{Agent Handles (3 conditions each)}},
  ] {};
\end{scope}

\end{tikzpicture}

  \caption{Testbed architecture. Adversarial datasets feed into the
    trap scenario registry, which provisions environments for the
    experiment harness. The harness deploys agents under test (one of
    four models in one of three conditions) into each environment,
    collects observations, computes metrics, and generates \LaTeX\
    assets via the reporter.}
  \label{fig:architecture}
\end{figure}

\subsection{Trap Scenario Interface}
\label{sec:testbed-interface}

Every trap scenario implements four methods:

\begin{itemize}[nosep,leftmargin=*]
  \item \textbf{\texttt{setup(config)}}: Provisions the adversarial
    environment---creates poisoned web pages, API endpoints,
    documents, or A2A messages. Returns a \texttt{TrapEnvironment}
    containing adversarial resources and ground-truth labels.

  \item \textbf{\texttt{execute(env, agent)}}: Deploys the agent
    under test into the environment and records its behavior as a
    \texttt{TrapObservation}, including whether the trap was
    triggered, detected, and escaped.

  \item \textbf{\texttt{evaluate(obs)}}: Computes six core metrics
    (\cref{sec:testbed-metrics}) from the observation.

  \item \textbf{\texttt{teardown(env)}}: Cleans up provisioned
    resources.
\end{itemize}

\subsection{The 22 Adversarial Scenarios}
\label{sec:testbed-scenarios}

\Cref{tab:scenarios} summarizes all 22 scenarios organized by
category. Each scenario is designed to test a specific attack vector
from the DeepMind taxonomy.

\begin{table}[t]
  \centering
  \caption{The 22 adversarial scenarios organized by DeepMind taxonomy
    category. Each row lists the scenario identifier, attack technique,
    and the primary metric used to assess attack success.}
  \label{tab:scenarios}
  \small
  \begin{tabular}{@{}llll@{}}
    \toprule
    \textbf{Cat.} & \textbf{ID} & \textbf{Technique} & \textbf{Metric} \\
    \midrule
    \multirow{4}{*}{\rotatebox[origin=c]{90}{\scriptsize Content Inj.}}
      & CI-1 & CSS invisible text  & Exec.\ rate \\
      & CI-2 & HTML comments        & Exec.\ rate \\
      & CI-3 & Image EXIF metadata  & Exec.\ rate \\
      & CI-4 & Dynamic cloaking     & Exec.\ rate \\
    \midrule
    \multirow{4}{*}{\rotatebox[origin=c]{90}{\scriptsize Semantic Man.}}
      & SM-1 & Authority framing    & Dec.\ drift \\
      & SM-2 & Emotional urgency    & Dec.\ drift \\
      & SM-3 & Context flooding     & Dec.\ drift \\
      & SM-4 & Identity manipulation & Dec.\ drift \\
    \midrule
    \multirow{4}{*}{\rotatebox[origin=c]{90}{\scriptsize Cognitive St.}}
      & CS-1 & Vector poisoning     & Accuracy $\downarrow$ \\
      & CS-2 & Ranking manipulation & Accuracy $\downarrow$ \\
      & CS-3 & Gradual drift        & Accuracy $\downarrow$ \\
      & CS-4 & Cross-contamination  & Accuracy $\downarrow$ \\
    \midrule
    \multirow{4}{*}{\rotatebox[origin=c]{90}{\scriptsize Behav.\ Ctrl.}}
      & BC-1 & Deceptive dialogs    & Wrong-act.\ rate \\
      & BC-2 & Misleading forms     & Wrong-act.\ rate \\
      & BC-3 & Hidden fields        & Wrong-act.\ rate \\
      & BC-4 & Infinite loops       & Token waste \\
    \midrule
    \multirow{3}{*}{\rotatebox[origin=c]{90}{\scriptsize Systemic}}
      & SY-1 & Message poisoning    & Cascade depth \\
      & SY-2 & Agent impersonation  & Blast radius \\
      & SY-3 & Cascade failure      & Cascade depth \\
    \midrule
    \multirow{3}{*}{\rotatebox[origin=c]{90}{\scriptsize HitL}}
      & HL-1 & Cherry-picked evidence & Manip.\ rate \\
      & HL-2 & Anchoring effects    & Manip.\ rate \\
      & HL-3 & Decision fatigue     & Manip.\ rate \\
    \bottomrule
  \end{tabular}
\end{table}

\subsection{Metrics Framework}
\label{sec:testbed-metrics}

We define six core metrics computed for every experiment cell:

\begin{itemize}[nosep,leftmargin=*]
  \item \metric{Trap Success Rate} ($\tau$): Binary indicator---1 if
    the agent executed the adversarial instruction, 0 otherwise.
  \item \metric{Detection Rate} ($\delta$): 1 if the agent flagged
    suspicious content, 0 otherwise.
  \item \metric{Escape Rate} ($\epsilon$): 1 if the agent both
    detected and avoided the trap.
  \item \metric{Decision Drift} ($\Delta$): Cosine distance between
    the agent's action and the ground-truth correct action.
  \item \metric{Time to Compromise} ($t_c$): Milliseconds from
    environment deployment to trap execution.
  \item \metric{Token Waste} ($w$): Completion tokens consumed
    during failed or manipulated processing.
\end{itemize}

For systemic scenarios, we additionally compute \metric{Cascade Depth}
(number of agents compromised in sequence) and \metric{Blast Radius}
(fraction of the agent crew affected).

\subsection{Experiment Matrix Visualization}
\label{sec:testbed-matrix}

\Cref{fig:matrix-flow} visualizes the full experiment matrix
design. The matrix generates 9{,}120 total experiment cells across
main, compound, and ablation sub-matrices.

\begin{figure}[t]
  \centering
  % ===========================================================================
% Experiment Matrix Flow Diagram
% ===========================================================================
% Included by: % ===========================================================================
% Experiment Matrix Flow Diagram
% ===========================================================================
% Included by: % ===========================================================================
% Experiment Matrix Flow Diagram
% ===========================================================================
% Included by: \input{figures/matrix-flow}
% Shows: the 22 × 4 × 3 × 10 experimental design
%
\begin{tikzpicture}[
  box/.style={
    rectangle,
    draw=black!70,
    fill=#1,
    rounded corners=2pt,
    minimum height=0.7cm,
    font=\small\sffamily,
    align=center,
    thick,
    inner sep=4pt,
  },
  box/.default=blue!8,
  bigbox/.style={
    rectangle,
    draw=black!50,
    fill=#1,
    rounded corners=4pt,
    minimum height=1.0cm,
    font=\sffamily,
    align=center,
    thick,
    inner sep=6pt,
  },
  bigbox/.default=white,
  arrow/.style={-Stealth, thick, black!70},
  multiply/.style={
    circle,
    draw=black!60,
    fill=gray!10,
    font=\small\bfseries,
    inner sep=2pt,
    minimum size=0.5cm,
  },
  total/.style={
    rectangle,
    draw=black!80,
    fill=red!10,
    rounded corners=3pt,
    minimum height=0.8cm,
    font=\small\sffamily\bfseries,
    align=center,
    thick,
    inner sep=5pt,
  },
  sublabel/.style={
    font=\scriptsize\sffamily,
    text=black!60,
  },
]

% ── Column 1: Scenarios ──
\node[box=blue!10] (scenarios) at (0, 0) {22 Scenarios};
\node[sublabel, below=1pt of scenarios] {6 categories};

% ── Multiply ──
\node[multiply] (x1) at (1.8, 0) {$\times$};

% ── Column 2: Models ──
\node[box=green!10] (models) at (3.4, 0) {4 Models};
\node[sublabel, below=1pt of models] {3 providers};

% ── Multiply ──
\node[multiply] (x2) at (5.0, 0) {$\times$};

% ── Column 3: Conditions ──
\node[box=orange!10] (conditions) at (6.8, 0) {3 Conditions};
\node[sublabel, below=1pt of conditions] {B / H / A};

% ── Multiply ──
\node[multiply] (x3) at (8.5, 0) {$\times$};

% ── Column 4: Repetitions ──
\node[box=purple!10] (reps) at (10.2, 0) {10 Reps};
\node[sublabel, below=1pt of reps] {statistical power};

% ── Arrows ──
\draw[arrow] (scenarios) -- (x1);
\draw[arrow] (x1) -- (models);
\draw[arrow] (models) -- (x2);
\draw[arrow] (x2) -- (conditions);
\draw[arrow] (conditions) -- (x3);
\draw[arrow] (x3) -- (reps);

% ── Sub-matrices (below) ──
\node[bigbox=blue!5] (main) at (1.7, -2.2) {
  \textbf{Main Matrix}\\
  {\small $22 \times 4 \times 2 \times 10$}\\
  {\small = 1,760 cells}
};

\node[bigbox=orange!5] (ablation) at (5.2, -2.2) {
  \textbf{Ablation Matrix}\\
  {\small $22 \times 4 \times 7 \times 10$}\\
  {\small = 6,160 cells}
};

\node[bigbox=green!5] (compound) at (9.0, -2.2) {
  \textbf{Compound Matrix}\\
  {\small $15 \times 4 \times 2 \times 10$}\\
  {\small = 1,200 cells}
};

% ── Convergence arrow ──
\node[total] (total) at (5.2, -4.2) {Total: 9,120 experiment cells};

\draw[arrow] (main.south) -- (total.north west);
\draw[arrow] (ablation.south) -- (total.north);
\draw[arrow] (compound.south) -- (total.north east);

% ── Top connection to sub-matrices ──
\draw[arrow, dashed, black!40] (reps.south) -- ++(0, -0.5) -| (main.north);
\draw[arrow, dashed, black!40] (reps.south) -- ++(0, -0.5) -| (ablation.north);
\draw[arrow, dashed, black!40] (reps.south) -- ++(0, -0.5) -| (compound.north);

\end{tikzpicture}

% Shows: the 22 × 4 × 3 × 10 experimental design
%
\begin{tikzpicture}[
  box/.style={
    rectangle,
    draw=black!70,
    fill=#1,
    rounded corners=2pt,
    minimum height=0.7cm,
    font=\small\sffamily,
    align=center,
    thick,
    inner sep=4pt,
  },
  box/.default=blue!8,
  bigbox/.style={
    rectangle,
    draw=black!50,
    fill=#1,
    rounded corners=4pt,
    minimum height=1.0cm,
    font=\sffamily,
    align=center,
    thick,
    inner sep=6pt,
  },
  bigbox/.default=white,
  arrow/.style={-Stealth, thick, black!70},
  multiply/.style={
    circle,
    draw=black!60,
    fill=gray!10,
    font=\small\bfseries,
    inner sep=2pt,
    minimum size=0.5cm,
  },
  total/.style={
    rectangle,
    draw=black!80,
    fill=red!10,
    rounded corners=3pt,
    minimum height=0.8cm,
    font=\small\sffamily\bfseries,
    align=center,
    thick,
    inner sep=5pt,
  },
  sublabel/.style={
    font=\scriptsize\sffamily,
    text=black!60,
  },
]

% ── Column 1: Scenarios ──
\node[box=blue!10] (scenarios) at (0, 0) {22 Scenarios};
\node[sublabel, below=1pt of scenarios] {6 categories};

% ── Multiply ──
\node[multiply] (x1) at (1.8, 0) {$\times$};

% ── Column 2: Models ──
\node[box=green!10] (models) at (3.4, 0) {4 Models};
\node[sublabel, below=1pt of models] {3 providers};

% ── Multiply ──
\node[multiply] (x2) at (5.0, 0) {$\times$};

% ── Column 3: Conditions ──
\node[box=orange!10] (conditions) at (6.8, 0) {3 Conditions};
\node[sublabel, below=1pt of conditions] {B / H / A};

% ── Multiply ──
\node[multiply] (x3) at (8.5, 0) {$\times$};

% ── Column 4: Repetitions ──
\node[box=purple!10] (reps) at (10.2, 0) {10 Reps};
\node[sublabel, below=1pt of reps] {statistical power};

% ── Arrows ──
\draw[arrow] (scenarios) -- (x1);
\draw[arrow] (x1) -- (models);
\draw[arrow] (models) -- (x2);
\draw[arrow] (x2) -- (conditions);
\draw[arrow] (conditions) -- (x3);
\draw[arrow] (x3) -- (reps);

% ── Sub-matrices (below) ──
\node[bigbox=blue!5] (main) at (1.7, -2.2) {
  \textbf{Main Matrix}\\
  {\small $22 \times 4 \times 2 \times 10$}\\
  {\small = 1,760 cells}
};

\node[bigbox=orange!5] (ablation) at (5.2, -2.2) {
  \textbf{Ablation Matrix}\\
  {\small $22 \times 4 \times 7 \times 10$}\\
  {\small = 6,160 cells}
};

\node[bigbox=green!5] (compound) at (9.0, -2.2) {
  \textbf{Compound Matrix}\\
  {\small $15 \times 4 \times 2 \times 10$}\\
  {\small = 1,200 cells}
};

% ── Convergence arrow ──
\node[total] (total) at (5.2, -4.2) {Total: 9,120 experiment cells};

\draw[arrow] (main.south) -- (total.north west);
\draw[arrow] (ablation.south) -- (total.north);
\draw[arrow] (compound.south) -- (total.north east);

% ── Top connection to sub-matrices ──
\draw[arrow, dashed, black!40] (reps.south) -- ++(0, -0.5) -| (main.north);
\draw[arrow, dashed, black!40] (reps.south) -- ++(0, -0.5) -| (ablation.north);
\draw[arrow, dashed, black!40] (reps.south) -- ++(0, -0.5) -| (compound.north);

\end{tikzpicture}

% Shows: the 22 × 4 × 3 × 10 experimental design
%
\begin{tikzpicture}[
  box/.style={
    rectangle,
    draw=black!70,
    fill=#1,
    rounded corners=2pt,
    minimum height=0.7cm,
    font=\small\sffamily,
    align=center,
    thick,
    inner sep=4pt,
  },
  box/.default=blue!8,
  bigbox/.style={
    rectangle,
    draw=black!50,
    fill=#1,
    rounded corners=4pt,
    minimum height=1.0cm,
    font=\sffamily,
    align=center,
    thick,
    inner sep=6pt,
  },
  bigbox/.default=white,
  arrow/.style={-Stealth, thick, black!70},
  multiply/.style={
    circle,
    draw=black!60,
    fill=gray!10,
    font=\small\bfseries,
    inner sep=2pt,
    minimum size=0.5cm,
  },
  total/.style={
    rectangle,
    draw=black!80,
    fill=red!10,
    rounded corners=3pt,
    minimum height=0.8cm,
    font=\small\sffamily\bfseries,
    align=center,
    thick,
    inner sep=5pt,
  },
  sublabel/.style={
    font=\scriptsize\sffamily,
    text=black!60,
  },
]

% ── Column 1: Scenarios ──
\node[box=blue!10] (scenarios) at (0, 0) {22 Scenarios};
\node[sublabel, below=1pt of scenarios] {6 categories};

% ── Multiply ──
\node[multiply] (x1) at (1.8, 0) {$\times$};

% ── Column 2: Models ──
\node[box=green!10] (models) at (3.4, 0) {4 Models};
\node[sublabel, below=1pt of models] {3 providers};

% ── Multiply ──
\node[multiply] (x2) at (5.0, 0) {$\times$};

% ── Column 3: Conditions ──
\node[box=orange!10] (conditions) at (6.8, 0) {3 Conditions};
\node[sublabel, below=1pt of conditions] {B / H / A};

% ── Multiply ──
\node[multiply] (x3) at (8.5, 0) {$\times$};

% ── Column 4: Repetitions ──
\node[box=purple!10] (reps) at (10.2, 0) {10 Reps};
\node[sublabel, below=1pt of reps] {statistical power};

% ── Arrows ──
\draw[arrow] (scenarios) -- (x1);
\draw[arrow] (x1) -- (models);
\draw[arrow] (models) -- (x2);
\draw[arrow] (x2) -- (conditions);
\draw[arrow] (conditions) -- (x3);
\draw[arrow] (x3) -- (reps);

% ── Sub-matrices (below) ──
\node[bigbox=blue!5] (main) at (1.7, -2.2) {
  \textbf{Main Matrix}\\
  {\small $22 \times 4 \times 2 \times 10$}\\
  {\small = 1,760 cells}
};

\node[bigbox=orange!5] (ablation) at (5.2, -2.2) {
  \textbf{Ablation Matrix}\\
  {\small $22 \times 4 \times 7 \times 10$}\\
  {\small = 6,160 cells}
};

\node[bigbox=green!5] (compound) at (9.0, -2.2) {
  \textbf{Compound Matrix}\\
  {\small $15 \times 4 \times 2 \times 10$}\\
  {\small = 1,200 cells}
};

% ── Convergence arrow ──
\node[total] (total) at (5.2, -4.2) {Total: 9,120 experiment cells};

\draw[arrow] (main.south) -- (total.north west);
\draw[arrow] (ablation.south) -- (total.north);
\draw[arrow] (compound.south) -- (total.north east);

% ── Top connection to sub-matrices ──
\draw[arrow, dashed, black!40] (reps.south) -- ++(0, -0.5) -| (main.north);
\draw[arrow, dashed, black!40] (reps.south) -- ++(0, -0.5) -| (ablation.north);
\draw[arrow, dashed, black!40] (reps.south) -- ++(0, -0.5) -| (compound.north);

\end{tikzpicture}

  \caption{Experiment matrix flow. 22~scenarios are crossed with
    4~models and 3~conditions (baseline, hardened, ablated) at
    10~repetitions each. Compound and ablation sub-matrices expand
    the total to 9{,}120 experiment cells.}
  \label{fig:matrix-flow}
\end{figure}
